Na camada de segurança, é de autoria da equipe a biblioteca \textit{AsconSecure}, descrita na Seção~\ref{subsec:asconsecure}, com o provisionamento da chave no próprio nó, o contador persistente que garante a unicidade do \textit{nonce} e o registro contra repetição no receptor. São também da equipe os programas de teste e de medição que produziram os resultados da Seção~\ref{sec:resultados_ascon}, que são a bateria de 4.907 adulterações de pacote, a suíte que mede memória, tempo de execução e tempo constante, o comparativo com o AES-128-GCM e a medição de energia por rajadas, além da análise desses resultados. Não são de autoria da equipe a implementação do Ascon-128a, que vem da biblioteca \textit{ascon-suite} \cite{weatherley2023ascon}, os vetores oficiais de teste, distribuídos com ela, e a implementação do AES-128-GCM, que vem da \textit{mbedTLS} incluída no ambiente de desenvolvimento do ESP32.

Como resultado, o trabalho contribui com a comparação entre o Ascon-128a em \textit{software} e o AES-128-GCM com acelerador de \textit{hardware} em dois microcontroladores da família ESP32, em tempo, memória e energia, e com a decomposição do custo do AES-GCM entre a cifra e a autenticação. Contribui também com a constatação de que, no ESP32-S3, o caminho por DMA impõe um custo fixo por chamada que domina o tempo das mensagens pequenas. O código e os registros das medições estão no repositório do projeto.
